%% tools/bench.tex -- the document tools/bench.sh times for the margin: ten pages of plain
%% text in the grown margin, the last one short so that the garden also grows under it.
%% The same file for every release (it uses nothing a release added), so their times compare.
%%   lualatex "\def\florilegiumseed{7}\def\benchoptions{}%% tools/bench.tex -- the document tools/bench.sh times for the margin: ten pages of plain
%% text in the grown margin, the last one short so that the garden also grows under it.
%% The same file for every release (it uses nothing a release added), so their times compare.
%%   lualatex "\def\florilegiumseed{7}\def\benchoptions{}%% tools/bench.tex -- the document tools/bench.sh times for the margin: ten pages of plain
%% text in the grown margin, the last one short so that the garden also grows under it.
%% The same file for every release (it uses nothing a release added), so their times compare.
%%   lualatex "\def\florilegiumseed{7}\def\benchoptions{}%% tools/bench.tex -- the document tools/bench.sh times for the margin: ten pages of plain
%% text in the grown margin, the last one short so that the garden also grows under it.
%% The same file for every release (it uses nothing a release added), so their times compare.
%%   lualatex "\def\florilegiumseed{7}\def\benchoptions{}\input{bench.tex}"
%% \benchoptions: more package options after grown (",nomargin" times the text alone).
\documentclass[11pt,a4paper]{article}
\renewcommand\rmdefault{cmr}\renewcommand\sfdefault{cmss}\renewcommand\ttdefault{cmtt}
\usepackage[T1]{fontenc}
\usepackage[a4paper,left=68mm,right=22mm,top=32mm,bottom=30mm,headsep=8mm]{geometry}
\providecommand\benchoptions{}
\edef\benchload{\noexpand\usepackage[grown\benchoptions]{florilegium}}\benchload
\pagestyle{plain}
\newcount\benchpar
\newcommand\benchtext{A florilegium was a gathering of flowers: passages copied out of
  longer books and bound together, so that a reader could carry the best of a library in one
  volume. This page holds no such passages. It is plain text set in a column, a paragraph
  repeated until the pages are full, so that the time spent on it is the time the margin
  takes and the time TeX takes to set the words, and nothing else. Each page is grown from
  the same seed whichever release draws it, and the last page ends early so that the garden
  also spreads into the space under its final lines.\par}
\begin{document}
\benchpar=0
\loop\benchtext\advance\benchpar 1 \ifnum\benchpar<57 \repeat
\florilegiumtextend
\end{document}
"
%% \benchoptions: more package options after grown (",nomargin" times the text alone).
\documentclass[11pt,a4paper]{article}
\renewcommand\rmdefault{cmr}\renewcommand\sfdefault{cmss}\renewcommand\ttdefault{cmtt}
\usepackage[T1]{fontenc}
\usepackage[a4paper,left=68mm,right=22mm,top=32mm,bottom=30mm,headsep=8mm]{geometry}
\providecommand\benchoptions{}
\edef\benchload{\noexpand\usepackage[grown\benchoptions]{florilegium}}\benchload
\pagestyle{plain}
\newcount\benchpar
\newcommand\benchtext{A florilegium was a gathering of flowers: passages copied out of
  longer books and bound together, so that a reader could carry the best of a library in one
  volume. This page holds no such passages. It is plain text set in a column, a paragraph
  repeated until the pages are full, so that the time spent on it is the time the margin
  takes and the time TeX takes to set the words, and nothing else. Each page is grown from
  the same seed whichever release draws it, and the last page ends early so that the garden
  also spreads into the space under its final lines.\par}
\begin{document}
\benchpar=0
\loop\benchtext\advance\benchpar 1 \ifnum\benchpar<57 \repeat
\florilegiumtextend
\end{document}
"
%% \benchoptions: more package options after grown (",nomargin" times the text alone).
\documentclass[11pt,a4paper]{article}
\renewcommand\rmdefault{cmr}\renewcommand\sfdefault{cmss}\renewcommand\ttdefault{cmtt}
\usepackage[T1]{fontenc}
\usepackage[a4paper,left=68mm,right=22mm,top=32mm,bottom=30mm,headsep=8mm]{geometry}
\providecommand\benchoptions{}
\edef\benchload{\noexpand\usepackage[grown\benchoptions]{florilegium}}\benchload
\pagestyle{plain}
\newcount\benchpar
\newcommand\benchtext{A florilegium was a gathering of flowers: passages copied out of
  longer books and bound together, so that a reader could carry the best of a library in one
  volume. This page holds no such passages. It is plain text set in a column, a paragraph
  repeated until the pages are full, so that the time spent on it is the time the margin
  takes and the time TeX takes to set the words, and nothing else. Each page is grown from
  the same seed whichever release draws it, and the last page ends early so that the garden
  also spreads into the space under its final lines.\par}
\begin{document}
\benchpar=0
\loop\benchtext\advance\benchpar 1 \ifnum\benchpar<57 \repeat
\florilegiumtextend
\end{document}
"
%% \benchoptions: more package options after grown (",nomargin" times the text alone).
\documentclass[11pt,a4paper]{article}
\renewcommand\rmdefault{cmr}\renewcommand\sfdefault{cmss}\renewcommand\ttdefault{cmtt}
\usepackage[T1]{fontenc}
\usepackage[a4paper,left=68mm,right=22mm,top=32mm,bottom=30mm,headsep=8mm]{geometry}
\providecommand\benchoptions{}
\edef\benchload{\noexpand\usepackage[grown\benchoptions]{florilegium}}\benchload
\pagestyle{plain}
\newcount\benchpar
\newcommand\benchtext{A florilegium was a gathering of flowers: passages copied out of
  longer books and bound together, so that a reader could carry the best of a library in one
  volume. This page holds no such passages. It is plain text set in a column, a paragraph
  repeated until the pages are full, so that the time spent on it is the time the margin
  takes and the time TeX takes to set the words, and nothing else. Each page is grown from
  the same seed whichever release draws it, and the last page ends early so that the garden
  also spreads into the space under its final lines.\par}
\begin{document}
\benchpar=0
\loop\benchtext\advance\benchpar 1 \ifnum\benchpar<57 \repeat
\florilegiumtextend
\end{document}
